\subsection{Propiedades de descenso, de paso al límite y de constructibilidad}
\label{subsection:IV.17.7}
\label{IV.17.7}

\begin{proposition}[17.7.1]
\label{IV.17.7.1}
Sean $f:X\to Y$ un morfismo localmente de presentación finita, $g:Y'\to Y$ un morfismo, $X'=X\times_Y Y'$, $f'=f_{(Y')}:X'\to Y'$ y $g':X'\to X$ las proyecciones canónicas. Sea $x'$ un punto de $X'$ y pongamos $x=g'(x')$, $y'=f'(x')$.
\begin{enumerate}
  \item[(i)] Si $f'$ es no ramificado en el punto $x'$, entonces $f$ es no ramificado en el punto $x$.
  \item[(ii)] Supongamos además que $g$ sea plano en el punto $y'$. Entonces, si $f'$ es liso (resp. étale) en el punto $x'$, $f$ es liso (resp. étale) en el punto $x$.
\end{enumerate}
\end{proposition}

Pongamos $y=f(x)=g(y')$, de modo que se tiene $f'^{-1}(y')=f^{-1}(y)\otimes_{k(y)}k(y')$; observemos que $f'$ es localmente de presentación finita.

(i) Como la propiedad de un morfismo (localmente de presentación finita) de ser no ramificado en un punto solo depende de la fibra del morfismo en dicho punto (17.4.1, $d$), podemos limitarnos al caso en que $Y$ e $Y'$ son espectros de cuerpos. Pero entonces equivale decir que $f$ (resp. $f'$) es no ramificado en $x$ (resp. $x'$) o que es étale en ese punto (17.6.1, $c$), por lo que (i) es consecuencia de (ii).

(ii) Como $f'$ es plano en el punto $x'$ (17.5.1), la hipótesis de que $g$ es plano en el punto $y'$ implica que $f$ es plano en el punto $x$, pues la proyección $g':X'\to X$ es un morfismo plano en el punto $x'$, y $f\circ g'=g\circ f'$ es un morfismo plano en el punto $X'$, de donde se sigue la conclusión (2.2.11, (iv)). El hecho de que $f$ sea liso (resp. étale) en el punto $x$ solo depende ya de la fibra $f^{-1}(y)$ (17.5.1 y 17.6.1), y nos reducimos, por tanto, una vez más al caso en que $Y=\Spec(k)$ e $Y'=\Spec(k')$ son espectros de cuerpos. Decir que $f'$ es liso en el punto $x'$ significa entonces (17.5.1) que $X'$ es un $k'$-preesquema geométricamente regular en el punto $x'$, y esto implica (6.7.8) que $X$ es un $k$-preesquema geométricamente regular en el punto $x$, y por tanto que $f$ es liso en el punto $x$. Supongamos además que $f'$ sea étale en el punto $x'$, de modo que $x'$ es aislado en $f'^{-1}(y')$ (17.6.1); como la proyección $f'^{-1}(y')\to f^{-1}(y)$ es un morfismo abierto (2.4.10), $x$ es aislado en $f^{-1}(y)$; como ya sabemos que $f$ es liso en el punto $x$, es étale en dicho punto (17.6.1).

\begin{corollary}[17.7.2]
\label{IV.17.7.2}
\begin{enumerate}
  \item[(i)] Con las notaciones de (17.7.1), sea $U$ (resp. $U'$) el conjunto de puntos de $X$ (resp. $X'$) donde $f$ (resp. $f'$) es no ramificado; entonces se tiene $U'=g'^{-1}(U)$.
  \item[(ii)] Supongamos además que $g$ sea plano, y sea $V$ (resp. $V'$) el conjunto de puntos de $X$ donde $f$ (resp. $f'$) es liso; entonces $V'=g'^{-1}(V)$.
\end{enumerate}
\end{corollary}

\begin{corollary}[17.7.3]
\label{IV.17.7.3}
\begin{enumerate}
  \item[(i)] Supongamos que $g$ es sobreyectivo; entonces, para que $f$ sea no ramificado, es necesario y suficiente que $f'$ lo sea.
  \item[(ii)] Sean $f:X\to Y$ un $S$-morfismo y $g:S'\to S$ un morfismo fielmente plano. Supongamos que $f$ es localmente de presentación finita, o que $g$ es cuasicompacto. Entonces, para que $f$ sea liso (resp. no ramificado, resp. étale), es necesario y suficiente que $f'$ lo sea.
\end{enumerate}
\end{corollary}

En (ii), el caso en que $f$ es localmente de presentación finita resulta de (17.7.1). Si $g$ es cuasicompacto y $f'$ es liso (resp. no ramificado, resp. étale), $f'$ es localmente de presentación finita por (2.7.1, (iv)) y nos reducimos al primer caso.

\begin{proposition}[17.7.4]
\label{IV.17.7.4}
Sean $f:X\to Y$ un morfismo, $g:Y'\to Y$ un morfismo plano y localmente de presentación finita, $X'=X\times_Y Y'$, $f'=f_{(Y')}:X'\to Y'$ y $g':X'\to X$ las

\oldpage[IV-4]{73}

proyecciones canónicas. Sea $V$ (resp. $V'$) el conjunto de los $x\in X$ (resp. $x'\in X'$) donde $f$ posee una de las propiedades siguientes (resp. donde $f'$ posee la misma propiedad): ser:
\begin{enumerate}
  \item[(i)] localmente de tipo finito;
  \item[(ii)] localmente de presentación finita;
  \item[(iii)] plano;
  \item[(iv)] no ramificado;
  \item[(v)] liso;
  \item[(vi)] étale.
\end{enumerate}
Entonces se tiene $V'=g'^{-1}(V)$ (dicho de otro modo, para que $f'$ tenga la propiedad considerada en un punto $x'$, es necesario y suficiente que $f$ tenga dicha propiedad en el punto $x$).
\end{proposition}

La propiedad (iii) solo se menciona como recordatorio, y no requiere la hipótesis de que $g$ sea localmente de presentación finita ((2.2.11, (iv)), teniendo en cuenta que la proyección $g':X'\to X$ es un morfismo plano). En virtud de (17.7.2), las afirmaciones relativas a las propiedades (iv), (v) y (vi) son consecuencias de la afirmación relativa a (ii). Basta, por tanto, considerar los casos (i) y (ii). Es claro que $V'\supset g'^{-1}(V)$; queda, pues, por demostrar que $V'\subset g'^{-1}(V)$; observemos que los conjuntos $V$ y $V'$ son abiertos, y $W=g'(V')$ también es abierto en $X$ en virtud de (2.4.6). Se trata, por tanto, de probar que el morfismo $f|W:W\to Y$ es localmente de tipo finito (resp. localmente de presentación finita); como, por hipótesis, el compuesto
\[
V'\xrightarrow{g''}W\xrightarrow{f|W}Y
\]
(donde $g''$ es la restricción de $g'$), igual a $V'\xrightarrow{f'}Y'\xrightarrow{g}Y$, es localmente de tipo finito (resp. localmente de presentación finita) y $g''$ es sobreyectivo (y, por tanto, fielmente plano), nos reducimos a probar el lema siguiente, que mejora (11.3.16):

\begin{lemma}[17.7.5]
\label{IV.17.7.5}
Sean $f:X\to Y$ un morfismo fielmente plano y localmente de presentación finita, $g:Y\to Z$ un morfismo tal que $g\circ f:X\to Z$ tenga una de las propiedades siguientes: ser:
\begin{enumerate}
  \item[(i)] localmente de tipo finito;
  \item[(ii)] localmente de presentación finita;
  \item[(iii)] de tipo finito.
\end{enumerate}
Entonces $g$ tiene la misma propiedad. Si, además, $f$ es cuasicompacto o $g$ cuasiseparado, la misma conclusión es válida para la propiedad:
\begin{enumerate}
  \item[(iv)] ser de presentación finita.
\end{enumerate}
\end{lemma}

En los casos (i) y (ii), se trata de ver que, para todo $y\in Y$, existen una vecindad abierta afín $V$ de $y$ en $Y$ y una vecindad abierta afín $W$ de $z=g(y)$ en $Z$ que contiene $g(V)$, tales que el morfismo $V\to W$, restricción de $g$, sea de tipo finito (resp. de presentación finita). Ahora bien, existen por hipótesis un $x\in X$ tal que $f(x)=y$, una vecindad abierta afín $U$ de $x$ en $X$, una vecindad abierta afín $V'$ de $y$ en $Y$ que contiene $f(U)$ y una vecindad abierta afín $W$ de $z$ en $Z$ que contiene $g(V')$, tales que el morfismo $f_1:U\to V'$, restricción de $f$, sea plano y de presentación finita, y el morfismo $g_1:V'\to W$, restricción de $g$, sea tal que $g_1\circ f_1$ sea de tipo finito (resp. de presentación finita). Entonces $f(U)$ es abierto en $Y$ (2.4.6), y si $V\subset f(U)$ es una vecindad afín de $y$, el morfismo $f_2:f_1^{-1}(V)\to V$,

\oldpage[IV-4]{74}

restricción de $f_1$, sigue siendo de presentación finita y es, además, fielmente plano; asimismo, si $g_2=g_1|V$, $g_2\circ f_2$ es de tipo finito (resp. de presentación finita), pues el abierto $f_1^{-1}(V)$ es cuasicompacto (resp. cuasicompacto y cuasiseparado). Nos reducimos entonces a las hipótesis de (11.3.16), de donde se concluye que $g_2$ es un morfismo de tipo finito (resp. de presentación finita).

En el caso (iii), la cuestión es local sobre $Z$, por lo que podemos suponer que $Z$ es afín, y resulta entonces que $X$ es cuasicompacto, y por tanto también $Y=f(X)$. El caso (iii) es, pues, consecuencia de (i).

En el caso (iv) (con las hipótesis adicionales sobre $f$ o $g$), también podemos suponer que $Z$ es afín, y por tanto que $X$ e $Y$ son cuasicompactos; además, ya sabemos que $g$ es localmente de presentación finita y cuasicompacto (1.1.3), de modo que todo se reduce a ver que $g$ es cuasiseparado, y basta, por tanto, mostrar que esta propiedad es cierta cuando se supone que $f$ es cuasicompacto. Ahora bien, puesto que $g\circ f$ es cuasiseparado, lo mismo ocurre con $f$ (1.2.2, (v)); como $f$ es cuasicompacto y localmente de presentación finita, es de presentación finita (1.6.1); basta entonces repetir el razonamiento del primer párrafo de la demostración de (11.3.16).

\begin{remark}[17.7.6]
\label{IV.17.7.6}
En la afirmación relativa a la propiedad (iv), no puede suprimirse la hipótesis de que $f$ es cuasicompacto; de lo contrario, ello implicaría que todo morfismo $g:Y\to Z$ cuasicompacto y localmente de presentación finita sería de presentación finita, conclusión que se sabe errónea (1.6.4). En efecto, podemos limitarnos al caso en que $Z$ es afín, y por tanto $Y$ cuasicompacto; existe entonces un recubrimiento finito $(U_i)$ de $Y$ por abiertos afines tales que las restricciones $g|U_i$ sean de presentación finita; bastaría entonces tomar como $X$ el preesquema \emph{suma} de los $U_i$, y como $f:X\to Y$ el morfismo canónico, que es evidentemente fielmente plano y localmente de presentación finita (1.4.3); $g\circ f$ sería de presentación finita en virtud de la elección de los $U_i$ y de (1.6.5), de donde se sigue nuestra afirmación.
\end{remark}

\begin{proposition}[17.7.7]
\label{IV.17.7.7}
Sean $f:Y\to S$ y $h:X\to S$ dos morfismos localmente de presentación finita, $g:X\to Y$ un $S$-morfismo, $x$ un punto de $X$, $y=g(x)$. Supongamos que $g$ sea plano en el punto $x$. Entonces, si $h$ es liso (resp. no ramificado, resp. étale) en el punto $x$, $f$ es liso (resp. no ramificado, resp. étale) en el punto $y$.
\end{proposition}

Pongamos $s=h(x)=f(y)$. Decir que $h$ es no ramificado en el punto $x$ (resp. que $f$ es no ramificado en el punto $y$) equivale a decir que $h^{-1}(s)$ es étale sobre $k(s)$ en el punto $x$ (resp. que $f^{-1}(s)$ es étale sobre $k(s)$ en el punto $y$) (17.4.1 y 17.6.1); como el morfismo $g_s:h^{-1}(s)\to f^{-1}(s)$ deducido de $g$ es plano en el punto $x$, vemos que podemos limitarnos a probar la proposición cuando $h$ es liso o étale en el punto $x$. Además, como $h$ es entonces plano en el punto $x$ (17.5.1), $f$ es plano en el punto $y$, como resulta de (2.2.11, (iv)). Equivale, por tanto, (17.5.1) decir que $f$ es liso (resp. étale) en el punto $y$, o que $f^{-1}(s)$ es liso (resp. étale) sobre $k(s)$ en el punto $x$. Nos reducimos así al caso en que $S=\Spec(k)$ es el espectro de un cuerpo.

(i) \emph{Caso de los morfismos lisos}. Como $g$ es un morfismo localmente de presentación finita (1.4.3, (v)), existe una vecindad abierta $U$ de $x$ en $X$ en la cual $g$ es plano (11.3.1) y $h$ liso. Además, $g(U)$ es una vecindad abierta de $y$ en $Y$ (2.4.6); reemplazando $X$ por $U$ e $Y$ por $g(U)$, podemos suponer, por tanto, que $g$ es \emph{fielmente plano} y que $h$ es liso, y nos reducimos a probar que $f$ es entonces liso. Si $k'$ es una extensión algebraicamente

\oldpage[IV-4]{75}

cerrada de $k$, $X\otimes_k k'$ es entonces liso sobre $k'$, y en virtud de (17.7.3, (ii)), basta probar que $Y\otimes_k k'$ es liso sobre $k'$ (puesto que $\Spec(k')\to\Spec(k)$ es fielmente plano); podemos, por tanto, limitarnos al caso en que $k$ es \emph{algebraicamente cerrado}. Como el conjunto de los puntos de $Y$ racionales sobre $k$ es entonces muy denso en $Y$ (10.4.8) y como el conjunto de los puntos de $Y$ donde $Y$ es liso sobre $k$ es abierto, vemos que basta probar que $Y$ es liso sobre $k$ en todo punto racional sobre $k$. Pero, en tal punto $y$, decir que $Y$ es liso sobre $k$ en dicho punto equivale a decir que $Y$ es \emph{regular} en $y$ (17.5.1 y 6.7.8). Ahora bien, se tiene $y=g(x)$ para un $x\in X$ y, por hipótesis, (17.5.1) $X$ es regular en el punto $x$; como $X$ e $Y$ son entonces localmente noetherianos y $g$ es plano, $Y$ es, en efecto, regular en el punto $y$ (6.5.1, (i)).

(ii) \emph{Caso de los morfismos étale}. Por (i), ya sabemos que $f$ es liso en el punto $y$; en virtud de (17.6.1), basta, por tanto, mostrar que $f$ es cuasifinito en el punto $y$, o también que $\sh{O}_{Y,y}$ es una $k$-álgebra finita. Ahora bien, como $\sh{O}_{X,x}$ es un $\sh{O}_{Y,y}$-módulo fielmente plano ($0_{\mathrm I}$, 6.6.2), $\sh{O}_{Y,y}$ se identifica con un $k$-submódulo de $\sh{O}_{X,x}$ y, como $\sh{O}_{X,x}$ es por hipótesis una $k$-álgebra finita, lo mismo ocurre con $\sh{O}_{Y,y}$.

\begin{proposition}[17.7.8]
\label{IV.17.7.8}
Con las notaciones de (8.8.1), se supone que $X_\alpha$ e $Y_\alpha$ son localmente de presentación finita sobre $S_\alpha$. Sean $f_\alpha:X_\alpha\to Y_\alpha$ un $S_\alpha$-morfismo y $f:X\to Y$ el $S$-morfismo correspondiente.
\begin{enumerate}
  \item[(i)] Sean $x$ un punto de $X$ y $x_\lambda$ su proyección canónica en $X_\lambda$. Para que $f$ sea liso (resp. no ramificado, resp. étale) en el punto $x$, es necesario y suficiente que exista $\lambda\geq\alpha$ tal que $f_\lambda$ sea liso (resp. no ramificado, resp. étale) en el punto $x_\lambda$.
  \item[(ii)] Supongamos además que $X_\alpha$ es cuasicompacto. Para que $f$ sea liso (resp. no ramificado, resp. étale), es necesario y suficiente que exista $\lambda\geq\alpha$ tal que $f_\lambda$ sea liso (resp. no ramificado, resp. étale).
\end{enumerate}
\end{proposition}

(i) Si $y=f(x)$, $y_\lambda=f_\lambda(x_\lambda)$ es la proyección canónica de $y$ en $Y_\lambda$, y se tiene $f^{-1}(y)=f_\lambda^{-1}(y_\lambda)\otimes_{k(y_\lambda)}k(y)$; la parte del enunciado relativa a los morfismos no ramificados resulta, por tanto, de (17.7.1, (i)), y basta, pues, considerar el caso de los morfismos lisos. Como $f$ y $f_\lambda$ son localmente de presentación finita, equivale decir que $f^{-1}(y)$ es geométricamente regular en el punto $x$ o que $f_\lambda^{-1}(y_\lambda)$ es geométricamente regular en el punto $x_\lambda$ (6.7.8). La proposición resulta, por tanto, de (17.5.1) y de (11.2.6).

(ii) Para todo $\lambda$, sea $U_\lambda$ el conjunto abierto de los $x_\lambda\in X_\lambda$ tales que $f_\lambda$ sea liso (resp. no ramificado, resp. étale) en el punto $x_\lambda$; sea $V_\lambda$ su imagen inversa en $X$. Como, por hipótesis, para todo $x\in X$ existe, en virtud de (i), un $\lambda$ tal que $f_\lambda$ sea liso (resp. no ramificado, resp. étale) en el punto $x_\lambda$, $X$ es unión de los $V_\lambda$. Por otra parte (17.3.3), para $\lambda\leq\mu$ se tiene $V_\lambda\subset V_\mu$, de modo que, como $X$ es cuasicompacto, existe un índice $\mu$ tal que $X=V_\mu$. Como los $X_\lambda$ son cuasicompactos, resulta entonces de (8.3.4) que existe un índice $\nu\geq\mu$ tal que la imagen inversa de $U_\mu$ en $X_\nu$ sea todo $X_\nu$, lo que significa que $f_\nu$ es liso (resp. no ramificado, resp. étale) por (17.3.3).

\begin{corollary}[17.7.9]
\label{IV.17.7.9}
Sean $S=\Spec(A)$ un esquema afín y $f:X\to S$ un morfismo. Las condiciones siguientes son equivalentes:
\begin{enumerate}
  \item[a)] $f$ es un morfismo de presentación finita y liso (resp. no ramificado, resp. étale).

  \oldpage[IV-4]{76}

  \item[b)] Existen un esquema afín noetheriano $S_0=\Spec(A_0)$, un morfismo de tipo finito $f_0:X_0\to S_0$ y un morfismo $S\to S_0$ tales que el $S$-preesquema $X_0\otimes_{S_0}S$ sea $S$-isomorfo a $X$ y que $f_0$ sea liso (resp. no ramificado, resp. étale).
  \item[c)] Se verifican las condiciones de b), y además $A_0$ es una $\mathbb{Z}$-subálgebra de tipo finito de $A$, correspondiendo el morfismo $S\to S_0$ a la inyección canónica $A_0\to A$.
\end{enumerate}
\end{corollary}

La demostración a partir de (17.7.8) es la misma que la de (11.2.7) a partir de (11.2.6).

\begin{proposition}[17.7.10]
\label{IV.17.7.10}
Sean $f:Y\to S$, $h:X\to S$ dos morfismos localmente de presentación finita, $g:X\to Y$ un $S$-morfismo, $x$ un punto de $X$, $y=g(x)$. Supongamos que $h$ sea plano en el punto $x$ y que $f$ sea no ramificado en el punto $y$. Entonces $f$ es étale en el punto $y$ y $g$ es plano en el punto $x$.
\end{proposition}

(Más adelante se verá (18.4.9) que, de hecho, se puede prescindir de la hipótesis de que $h$ es localmente de presentación finita).

Como la cuestión es local sobre $S$, $X$ e $Y$, podemos suponer que $S$, $X$ e $Y$ son afines, que $f,g,h$ son morfismos de presentación finita, que $h$ es plano y que $f$ es no ramificado. Teniendo en cuenta (11.2.7) y (17.7.9), podemos suponer además que $S$, $X$ e $Y$ son noetherianos; finalmente, podemos limitarnos al caso en que $S=\Spec(\sh{O}_{S,s})$ (donde $s=f(y)=h(x)$). Como el anillo $A=\sh{O}_{S,s}$ es un anillo local noetheriano, existen un anillo local noetheriano $B$ completo, de cuerpo residual algebraicamente cerrado, y un homomorfismo local $A\to B$ que hace de $B$ un $A$-módulo fielmente plano ($0_{\mathrm{III}}$, 10.3.1). Reemplazando $X$ e $Y$ por $X\times_S\Spec(B)$ e $Y\times_S\Spec(B)$, se concluye por (2.5.1) y (17.7.1) que podemos limitarnos al caso en que $A$ es completo y tiene un cuerpo residual algebraicamente cerrado. La hipótesis de que $f$ es no ramificado en el punto $y$ implica entonces (17.4.4) que $\sh{O}_{Y,y}$ es una $\sh{O}_{S,s}$-álgebra finita y un anillo local noetheriano completo ($0_{\mathrm I}$, 7.4.2), y como el cuerpo residual de $\sh{O}_{S,s}$ es algebraicamente cerrado, el homomorfismo $\sh{O}_{S,s}\to\sh{O}_{Y,y}$ es sobreyectivo (17.4.4). Pero, por otra parte, la hipótesis de que $h$ es plano en el punto $x$ implica que el homomorfismo compuesto $\sh{O}_{S,s}\to\sh{O}_{Y,y}\to\sh{O}_{X,x}$ es inyectivo ($0_{\mathrm I}$, 6.5.1), y por tanto el homomorfismo $\sh{O}_{S,s}\to\sh{O}_{Y,y}$ es biyectivo, lo que muestra que $f$ es étale en el punto $y$ (17.6.3); además, $\sh{O}_{X,x}$ es un $\sh{O}_{Y,y}$-módulo plano, y por tanto $g$ es plano en el punto $x$.

\begin{proposition}[17.7.11]
\label{IV.17.7.11}
Sea $S$ un preesquema, sean $X$ e $Y$ dos $S$-preesquemas localmente de presentación finita sobre $S$, y sea $f:X\to Y$ un $S$-morfismo. Para todo $s\in S$, se denotan por $X_s$, $Y_s$, $f_s$ los preesquemas y el morfismo deducidos de $X$, $Y$, $f$ por el cambio de base $\Spec(k(s))\to S$. Entonces:
\begin{enumerate}
  \item[(i)] El conjunto de los $x\in X$ tales que, si $s$ es la imagen de $x$ en $S$, $f_s:X_s\to Y_s$ sea liso (resp. no ramificado, resp. étale, resp. diferencialmente liso) en el punto $x$, es localmente constructible.
  \item[(ii)] Supongamos que $f$ es de presentación finita. El conjunto de los $y\in Y$ tales que, si $s$ es la imagen de $y$ en $S$, $f_s$ sea liso (resp. no ramificado, resp. étale, resp. diferencialmente liso) en todos los puntos de $f_s^{-1}(y)$, es localmente constructible.
  \item[(iii)] Supongamos que $X$ e $Y$ son de presentación finita sobre $S$. El conjunto de los $s\in S$ tales que $f_s$ sea liso (resp. no ramificado, resp. étale, resp. diferencialmente liso) es localmente constructible.
\end{enumerate}
\end{proposition}

Sea $E$ el conjunto de los $x\in X$ que verifican la propiedad considerada en (i). Entonces el conjunto $F$ de los $y\in Y$ que verifican la propiedad correspondiente considerada en (ii) no es

\oldpage[IV-4]{77}

otro que $Y-f(X-E)$, por lo que (ii) resulta de (i) y del teorema de Chevalley cuando $f$ es de presentación finita (1.8.4). De igual modo, si $h:Y\to S$ es el morfismo estructural, el conjunto de los $s\in S$ que verifican la propiedad correspondiente considerada en (iii) es $S-h(Y-F)$, por lo que (iii) se deduce también de (ii) y del teorema de Chevalley cuando $f$ y $h$ son de presentación finita. Basta, por tanto, probar las afirmaciones de (i).

Probemos primero (i) cuando se trata de la propiedad de ser \emph{liso}. Como la cuestión es local sobre $X$, podemos limitarnos al caso en que $S=\Spec(A)$, $X=\Spec(B)$ e $Y=\Spec(C)$ son afines, siendo $B$ y $C$ $A$-álgebras de presentación finita. Razonando como al principio de (9.9.1) y utilizando (17.7.2, (ii)), nos reducimos al caso en que $A$ es noetheriano. En virtud de ($0_{\mathrm{III}}$, 9.2.3), nos reducimos a ver que, si $x\in E$ (resp. si $x\notin E$), existe una vecindad $V$ de $x$ en $\overline{\{x\}}$ contenida en $E$ (resp. en $X-E$). Denotando por $g:X\to S$ y $h:Y\to S$ los morfismos estructurales, podemos primero reemplazar $S$ por el subpreesquema reducido $S'$ de $S$ que tiene $\overline{\{g(x)\}}$ por espacio subyacente, $X$ por $g^{-1}(S')$, e $Y$ por $h^{-1}(S')$, pues las fibras de $X$ y $X'$ (resp. $Y$ e $Y'$) en los puntos de $S'$ son las mismas. Dicho de otro modo, podemos limitarnos al caso en que $S$ es \emph{íntegro} y en que $\eta=g(x)=h(y)$ (donde $y=f(x)$) es su punto genérico.

$1^\circ$ Supongamos primero que $x\in E$. Los anillos locales $\sh{O}_{X_\eta,x}$ y $\sh{O}_{Y_\eta,y}$ son respectivamente iguales a $\sh{O}_{X,x}$ y $\sh{O}_{Y,y}$; como la propiedad de lisitud de un morfismo de presentación finita en un punto solo depende del anillo local de dicho punto y del anillo local de su imagen (17.5.1), vemos que la hipótesis $x\in E$ equivale a decir que el morfismo $f$ es liso en el punto $x$; conserva entonces esta propiedad en los puntos de una vecindad abierta de $x$ en $X$, y basta aplicar (17.3.3, (iii)) para obtener la conclusión.

$2^\circ$ Supongamos, en segundo lugar, que $x\in X-E$, y que el morfismo $f_\eta$ no sea plano en el punto $x$. La conclusión resulta entonces del lema siguiente, que precisa (11.2.8):

\begin{lemma}[17.7.11.1]
\label{IV.17.7.11.1}
Sean $g:X\to S$, $h:Y\to S$ dos morfismos localmente de presentación finita, $f:X\to Y$ un $S$-morfismo, y $\sh{F}$ un $\sh{O}_X$-módulo cuasicoherente de presentación finita. Entonces el conjunto $E$ de los $x\in X$ tales que $\sh{F}_{g(x)}$ sea $f_{g(x)}$-plano en el punto $x$ es localmente constructible.
\end{lemma}

Razonando una vez más como al principio de (9.9.1) y utilizando (2.5.1), nos reducimos al caso en que $S$, $X$ e $Y$ son noetherianos; después nos reducimos, como antes, al caso en que $S$ es íntegro, en que $\eta=g(x)=h(f(x))$ es su punto genérico, y se trata de mostrar que si $x\in E$ (resp. $x\notin E$) existe una vecindad $V$ de $x$ en $\overline{\{x\}}$ contenida en $E$ (resp. en $X-E$). El caso en que $x\in E$ es consecuencia inmediata de (11.1.1). Para tratar el caso en que $x\notin E$, se razona como en (9.4.7.1), cuyas notaciones conservamos, de modo que podemos suponer, reemplazando eventualmente $Y$ y $X$ por vecindades de $f(x)$ y $x$, respectivamente, que existen dos $\sh{O}_Y$-módulos coherentes $\sh{G}$, $\sh{H}$ y un $\sh{O}_Y$-homomorfismo $u:\sh{G}\to\sh{H}$ tales que, para todo $s\in S$, $u_s:\sh{G}_s\to\sh{H}_s$ sea inyectivo, pero que el homomorfismo $1\otimes u_\eta:\sh{F}_\eta\otimes_{\sh{O}_{Y_\eta}}\sh{G}_\eta\to\sh{F}_\eta\otimes_{\sh{O}_{Y_\eta}}\sh{H}_\eta$ no sea inyectivo en el punto $x$, dicho de otro modo, $x\in\operatorname{Supp}(\operatorname{Ker}(1\otimes u_\eta))$; si se pone $T=\overline{\{x\}}$, se tiene, por tanto, $\operatorname{Supp}(\operatorname{Ker}(1\otimes u_\eta))\supset T_\eta$. Pero podemos suponer que, para todo $s\in S$, se tiene $\operatorname{Ker}(1\otimes u_s)=(\operatorname{Ker}(1\otimes u))_s$ (9.4.2), y por tanto

\oldpage[IV-4]{78}

$\operatorname{Supp}(\operatorname{Ker}(1\otimes u_s))=(\operatorname{Supp}(\operatorname{Ker}(1\otimes u)))_s$ (I, 9.1.13.1); resulta finalmente de (9.5.2) que, para $s$ en una vecindad de $\eta$, se tiene $(\operatorname{Supp}(\operatorname{Ker}(1\otimes u)))_s\supset T_s$, lo que demuestra el lema.

$3^\circ$ Supongamos ahora que $x\in X-E$, que el morfismo $f_\eta$ sea plano en el punto $x$, pero que $f_\eta$ no sea liso en el punto $x$. Observemos que decir que $f_\eta$ es plano en el punto $x$ equivale a decir que el propio $f$ es plano en el punto $x$ y, reemplazando $X$ por una vecindad de $x$, podemos suponer que $f$ es plano (11.1.1); de ello se concluye que lo mismo ocurre con $f_s$ para todo $s\in S$, y como para todo $y\in Y$, $f^{-1}(y)=f_{h(y)}^{-1}(y)$, equivale decir que $f_{g(x')}$ es liso en el punto $x'$ o decir que $f$ es liso en el punto $x'$. Pero el conjunto de los $x'\in X$ donde $f$ es liso es \emph{abierto} en $X$ (12.1.7), por lo que el conjunto de los $x'\in X$ donde $f$ no es liso es cerrado, y puesto que contiene a $x$ por hipótesis, también contiene a $\overline{\{x\}}$, lo que termina la demostración para la primera propiedad considerada en (i).

Probemos, en segundo lugar, (i) cuando se trata de la propiedad de ser \emph{étale}. Observemos para ello que esta propiedad para $f_s$ en el punto $x$ equivale a decir que $f_s$ es a la vez liso y cuasifinito en el punto $x$ (17.6.1). Ahora bien, equivale decir que $f_{g(x)}$ es cuasifinito en el punto $x$ o que el propio $f$ es cuasifinito en dicho punto; resulta, por tanto, de (13.1.4) que el conjunto de los puntos $x$ tales que $f_{g(x)}$ sea cuasifinito en el punto $x$ es abierto en $X$, y \emph{a fortiori} localmente constructible; la conclusión resulta entonces de que el conjunto de los $x$ tales que $f_{g(x)}$ sea liso en el punto $x$ es también localmente constructible.

Pasemos a la demostración de (i) cuando se trata de la propiedad de ser \emph{diferencialmente liso}. Sea $p:X\times_Y X\to X$ la segunda proyección canónica; la segunda proyección $X_s\times_{Y_s}X_s\to X_s$ no es otra que $p_s$ para todo $s\in S$, y resulta, por tanto, de (17.12.5.1)\footnote{El lector comprobará que el resultado de (17.7.11) no se utiliza en lo que sigue del § 17 y que, por tanto, no hay círculo vicioso.} que para que $f_{g(x)}$ sea diferencialmente liso en el punto $x$, es necesario y suficiente que $p_{g(x)}$ sea liso en el punto $\Delta_f(x)$. Como $p$ es localmente de presentación finita, el conjunto de los puntos $z\in X\times_Y X$ tales que $p_{g(p(z))}$ sea liso en el punto $z$ es localmente constructible, y la intersección de este conjunto con el conjunto localmente cerrado $\Delta_f(X)$ es, por tanto, también localmente constructible en $\Delta_f(X)$ (1.8.2); como la restricción de $p$ a $\Delta_f(X)$ es un isomorfismo sobre $X$, vemos que el conjunto de los $x\in X$ tales que $f_{g(x)}$ sea diferencialmente liso en el punto $x$ es localmente constructible.

Consideremos, finalmente, la propiedad de ser no ramificado; observemos que el morfismo diagonal $\Delta_f:X\to X\times_Y X$ es una inmersión localmente de presentación finita (1.4.3.1) y, para todo $s\in S$, el morfismo diagonal $\Delta_{f_s}:X_s\to X_s\times_{Y_s}X_s$ no es otro que $(\Delta_f)_s$; decir que $f_{g(x)}$ es no ramificado en el punto $x$ equivale a decir que $(\Delta_f)_{g(x)}$ es un isomorfismo local en el punto $x$ (17.4.1), y puesto que es una inmersión localmente de presentación finita, equivale (17.9.1)\footnote{El lector comprobará que el resultado de (17.7.11) no se utiliza en el resto del § 17 y que, por tanto, no hay círculo vicioso.} a decir que $(\Delta_f)_{g(x)}$ es étale en el punto $x$; basta, por tanto, aplicar lo visto más arriba para la propiedad de ser étale.

\oldpage[IV-4]{79}
